% Official source: https://www.imo-official.org/assets/documents/problems/2019/2019_eng.pdf
\competitionyear{2019}

\begin{problem}{2019}{1}
Find all functions $f\colon\Z\to\Z$ such that
\[
    f(2a)+2f(b)=f(f(a+b))
\]
for all integers $a, b$.
\end{problem}

\begin{problem}{2019}{2}
In triangle $ABC$, let $A_1\in BC$ and $B_1\in AC$. Choose $P\in AA_1$ and $Q\in BB_1$ with $PQ\parallel AB$. Point $P_1$ lies on line $PB_1$, with $B_1$ strictly between $P$ and $P_1$, and $\angle PP_1C=\angle BAC$. Similarly, point $Q_1$ lies on line $QA_1$, with $A_1$ strictly between $Q$ and $Q_1$, and $\angle CQ_1Q=\angle CBA$.

Prove that $P, Q, P_1, Q_1$ are concyclic.
\end{problem}

\begin{problem}{2019}{3}
A social network has $2019$ users with mutual friendships. The following operation may be performed repeatedly: choose users $A, B, C$ such that $A$ is friends with both $B$ and $C$, while $B$ and $C$ are not friends. Make $B$ and $C$ friends and end the friendships $AB$ and $AC$, leaving every other friendship unchanged.

Initially, $1010$ users each have $1009$ friends, and the other $1009$ users each have $1010$ friends. Prove that some sequence of the allowed operations leaves every user with at most one friend.
\end{problem}

\begin{problem}{2019}{4}
Find all pairs $(k, n)$ of positive integers satisfying
\[
    k!=(2^n-1)(2^n-2)(2^n-4)\cdots(2^n-2^{n-1}).
\]
\end{problem}

\begin{problem}{2019}{5}
The Bank of Bath issues coins marked $H$ on one side and $T$ on the other. Harry arranges $n$ such coins in a row and repeats this operation: when exactly $k>0$ coins show $H$, he flips the $k$-th coin from the left; if no coins show $H$, he stops. For instance, starting with $THT$ for $n=3$ yields
\[
    THT\to HHT\to HTT\to TTT,
\]
so the process ends after three operations.
\begin{problemsubparts}
    \item Prove that the process terminates for every initial configuration.
    \item For a configuration $C$, let $L(C)$ count the operations before termination. Thus $L(THT)=3$ and $L(TTT)=0$. Determine the average of $L(C)$ over all $2^n$ initial configurations $C$.
\end{problemsubparts}
\end{problem}

\begin{problem}{2019}{6}
Let $I$ be the incentre of an acute triangle $ABC$ with $AB\ne AC$. Its incircle $\omega$ touches $BC$, $CA$, and $AB$ at $D$, $E$, and $F$, respectively. The line through $D$ perpendicular to $EF$ meets $\omega$ again at $R$, and $AR$ meets $\omega$ again at $P$. The circumcircles of triangles $PCE$ and $PBF$ meet again at $Q$.

Prove that $DI$ and $PQ$ intersect on the line through $A$ perpendicular to $AI$.
\end{problem}
